\documentclass[12pt]{article}
\usepackage[margin=1in]{geometry}
\usepackage{enumitem}
\usepackage{titlesec}
\usepackage{parskip}
\usepackage{graphicx}
\usepackage{xcolor}
\usepackage{hyperref}
\usepackage{fontenc}

\title{\textbf{Muller 1997 RG Notes}\\
\large Müller, Wolfgang C.. ``Inside the Black Box: A Confrontation of Party Executive Behaviour and Theories of Party Organizational Change'' [1997].\\
\textit{Party Politics} 3(3): 293-313.}
\author{}
\date{}

\begin{document}

\maketitle

Central idea: in early formations of mass parties (especially post-WW2), the party organization itself was the primary means of outreach, engagement, and campaign operations. Later, especially as these mass parties focused more attention towards voters who were not already directly affiliated, they instead turned towards surveys and mass media as primary mechanisms of engagement and two-way communication

As new mass parties became more mainstream (example given is SPÖ in Austria, who gained significant prestige over their first ten years of operations), it became clear that campaigning towards existing party members was inefficient: elections are decided by the ``on-the-fence'' cohorts, who will inherently have less direct engagement with the party institutions

Broader connection to Stokes et al --- they claim that party operatives (particularly brokers) are still a primary means of communications between voters and parties, which this article suggests is not always the case, especially for mid-century mass movements

Also connect to Diermeier and Vlaicu(2011) --- the party as an institution may be less necessary on the legislative side as well, particularly if interests are already aligned

\end{document}
